\documentclass[11pt,a4paper]{article}
\usepackage[T1]{fontenc}
\usepackage[utf8]{inputenc}
\usepackage{lmodern}
\usepackage[a4paper,margin=27mm,headheight=15pt]{geometry}
\usepackage{amsmath,amssymb,amsthm,mathtools,mathrsfs}
\usepackage{microtype}
\usepackage{booktabs,array,longtable}
\usepackage{xcolor}
\usepackage{xurl}
\usepackage{enumitem}
\usepackage{needspace}
\usepackage{fancyhdr}
\usepackage{hyperref}
\hypersetup{colorlinks=true,linkcolor=blue!45!black,citecolor=blue!45!black,
urlcolor=blue!45!black,
pdftitle={Finite Resonance Control of Nonlinear Hahn--Dulac Transseries},
pdfauthor={Research article prepared for Vladimir Reshetnikov},
pdfsubject={Exact logarithmic slopes, finite algebraic certificates, nonlinear differential equations, and accumulating exponents}}
\urlstyle{same}
\numberwithin{equation}{section}
\newtheorem{theorem}{Theorem}[section]
\newtheorem{proposition}[theorem]{Proposition}
\newtheorem{lemma}[theorem]{Lemma}
\newtheorem{corollary}[theorem]{Corollary}
\theoremstyle{definition}
\newtheorem{definition}[theorem]{Definition}
\newtheorem{example}[theorem]{Example}
\theoremstyle{remark}
\newtheorem{remark}[theorem]{Remark}
% ed. (2026-09-29): unnumbered environment for ProveIt editorial notes; remarks share the theorem counter, so the notes must not be numbered.
\newtheorem*{ednote}{Editorial note (ProveIt, 2026-09-29)}
\newcommand{\CC}{\mathbb C}
\newcommand{\RR}{\mathbb R}
\newcommand{\NN}{\mathbb N}
\newcommand{\G}{\Gamma}
\newcommand{\Reson}{\mathcal R}
\newcommand{\Div}{\operatorname{Div}}
\newcommand{\supp}{\operatorname{supp}}
\newcommand{\val}{\operatorname{v}}
\newcommand{\spec}{\operatorname{spec}}
\newcommand{\dd}{\partial_L}
\newcommand{\norm}[1]{\left\lVert #1\right\rVert}
\newcommand{\abs}[1]{\left\lvert #1\right\rvert}
\newcommand{\floor}[1]{\left\lfloor #1\right\rfloor}
\newcommand{\ceil}[1]{\left\lceil #1\right\rceil}
\newcommand{\Log}{\operatorname{Log}}
\newcommand{\snap}{\texttt{3d5973524506411392a911470b5ddc35521568ea}}
\setlength{\parskip}{3pt}
\setlength{\parindent}{1.2em}
\setlist{itemsep=3pt,topsep=5pt,leftmargin=1.8em}
\linespread{1.025}
\pagestyle{fancy}
\fancyhf{}
\fancyhead[L]{\small Nonlinear Hahn--Dulac transseries}
\fancyhead[R]{\small Finite resonance control}
\fancyfoot[C]{\thepage}
\renewcommand{\headrulewidth}{0.3pt}
\setcounter{tocdepth}{1}

\begin{document}
\hypersetup{pageanchor=false}
\begin{titlepage}
\thispagestyle{empty}
\vspace*{9mm}
\noindent{\large\scshape Transseries research}\par
\vspace{10mm}
\noindent{\LARGE\bfseries Finite Resonance Control of\par}
\vspace{3mm}
\noindent{\LARGE\bfseries Nonlinear Hahn--Dulac Transseries\par}
\vspace{7mm}
\noindent{\large Exact logarithmic slopes, analytic realization,\par
and accumulation thresholds}\par
\vspace{10mm}
\hrule
\vspace{6mm}
\noindent Research article prepared for Vladimir Reshetnikov\par
\noindent 29 September 2026
\vspace{8mm}
\begin{abstract}
We study the nonlinear regular-singular system
\( (x\,d/dx-A)y=F(x,y) \), where \(A\) has real spectrum and the
logarithm-free coefficients of \(F\) have support in a positive real
Hahn semigroup. The support may have finite accumulation points.
Every choice of the allowed resonant constants determines a unique formal
solution \(y=\sum_{\gamma>0}x^\gamma P_\gamma(\log x)\).
Although its logarithmic degrees can be unbounded, their maximal ratio
to the exponent is determined by finitely many resonant polynomial
blocks. The determining polynomial blocks depend polynomially on only finitely
many input coordinates. The slope sublevel sets are algebraic, and the
maximum is invariant under invertible
logarithm-free local changes of the dependent variables.

We also prove a universal analytic-realization criterion: every absolutely
convergent input has a convergent normalized Hahn--Dulac solution precisely
when the nonresonant semigroup exponents stay uniformly separated from the
eigenvalues of \(A\). A weighted polynomial norm supplies explicit
contraction and remainder certificates. Failure of the gap condition
admits arbitrarily small, logarithm-free counterexamples.

An explicit nonlinear Riccati family with actions \(2-1/n\) separates
finite resonance control from convergence particularly sharply. Its
logarithmic slope depends on the single polynomial \(b+ca^2\), whereas
absolute convergence is equivalent to \(\sum n|t_n|<\infty\). Thus
coefficients \(t_n=\varepsilon n^{-p}\) have the exact threshold \(p>2\),
even in the presence of nonlinear feedback. Complete conventional proofs,
finite symbolic verification, and twelve further research directions are
provided. Classical Dulac majorants and the repository's linear
resonance results are acknowledged; historical priority is not asserted.
\end{abstract}
\vfill
\noindent\small\textbf{Scope.} Positive real actions; one logarithm;
first-order systems with no derivatives inside \(F\). These are human-readable
proofs, not a Lean formalization or an independently peer-reviewed report.\par
\vspace{2mm}
\noindent\small\textbf{Inspected repository snapshot.}\quad\snap
\end{titlepage}
\hypersetup{pageanchor=true}
\tableofcontents
\newpage

\section{The problem, the contribution, and its boundaries}
\label{sec:introduction}

Two distinct kinds of finiteness occur in transseries calculations. A
coefficient can have only finitely many algebraic ancestors even when
infinitely many exponents lie below it. Conversely, finitely many
resonances can generate infinitely many nonzero logarithmic powers once
nonlinear multiplication is allowed. Confusing these two facts leads to
an inappropriate choice of ambient space and, more seriously, to
incorrect inferences about convergence.

This article asks a concrete extension question:
\begin{quote}
For a nonlinear regular-singular system with well-ordered real actions,
can all logarithmic growth be determined from finite resonance data,
and when does the resulting formal solution define an actual convergent
function?
\end{quote}
The answer has three parts. The correct global invariant is a
\emph{degree-to-action slope}, not a finite bound on logarithmic degree.
That slope is determined by a finite algebraic certificate. Convergence,
however, can depend on infinitely many coefficients that do not occur in
that certificate.

\subsection{Relation to the inspected ProveIt materials}
The repository's support module contains additive versions of Neumann's
support and finite-decomposition lemmas, obtained through explicit
Mathlib bridges \cite{repo-support}. The finite-divisor argument used
below is a direct application of that established infrastructure.
The residue-obstruction package treats finite obstructions and
logarithmic-depth promotion for linear operators
\cite{repo-residue}. The Hahn--Fuchsian package gives a finite resonance
matrix for a linear system, a logarithmic-degree filtration, and a sharp
universal convergence criterion for its normalized gauge
\cite{repo-fuchsian}.

We do not rebrand those linear results as nonlinear discoveries. Here
there is generally no finite-dimensional vector space of solutions and
no single nilpotent matrix whose nilpotency index bounds every
logarithmic degree. Indeed, the scalar equation in
Section~\ref{sec:examples} has coefficients of degree exactly \(n\) at
action \(n\). The replacement proved here is an exact finite-resonance
formula for the maximal slope, together with nonlinear algebraic
cancellation loci and invariance under local coordinate changes.
Our analytic gap measures distances to \emph{eigenvalues}, whereas the
linear gauge-normalization problem naturally involves
\emph{eigenvalue differences}. They are different operators.
% ed. (2026-09-29): an independent second proof of thm:universal in a package of the same batch, and related later packages
\begin{ednote}
A package of the same batch, written independently,
\path{Nonlinear_Hahn_Fuchsian_Algebraic_Convergence_Loci/}, proves the
universal criterion of Theorem~\ref{thm:universal} (positivity of the
distance from the positive exponents to the eigenvalues they do not equal)
for logarithm-free positive solutions and arbitrary complex \(A\) (its
\texttt{thm:gap}), with genuinely quadratic counterexamples: one criterion,
two independent proofs, which agree where both apply. That package also
shows that nonlinear products of separated forcing can create the small
divisors (its \texttt{thm:hidden}, again with the threshold \(p>2\)); the
example of Section~\ref{sec:accumulation} instead has accumulating forcing.
For linear gauges with prescribed edge supports the universal quantifier
over all of \(\G\) is refined path by path in
\path{Path_Sensitive_Small_Divisors_Hahn_Fuchsian/}. The exact logarithmic
degree of linear scalar equations is determined by Smith invariants in
\path{Exact_Logarithmic_Degree_Smith_Invariants/}, which was filed after
the snapshot this article read. Theorem~\ref{thm:strata} is a nonlinear
counterpart of the Hahn--Fuchsian article's question ``Geometry of the
logarithmic strata'' \cite{repo-fuchsian}, not an answer to it.
\end{ednote}

\subsection{Established ingredients and proposed additions}
General fixed-point methods for generalized series are established;
van der Hoeven proves an implicit-function theorem for Noetherian
operators \cite{vdh}. Our elementary valuation construction is a
special-purpose instance, not a new general implicit-function theorem.
Gontsov and Goryuchkina prove a majorant convergence theorem for
generalized power series \cite{gg-generalized}.

In their Dulac-series paper, Goryuchkina and Gontsov already obtain a
linear degree bound in the integer action and use coefficient norms,
rescaling of the logarithmic variable, and bounds for polynomial Euler
operators \cite[Lemmas 2--3]{gg-dulac}. Neither the existence of a linear
degree bound nor that norm method is claimed as new here. We adapt those
ingredients to finite resonances and possibly non-locally-finite Hahn
supports.

The proposed additions are the exact maximum formula of
Theorem~\ref{thm:slope}, its finite algebraic certificate and
coordinate-invariance refinements, the nonlinear formulation of the
universal gap equivalence, and the input-specific nonlinear accumulation
theorem of Section~\ref{sec:accumulation}. These statements are proved
below. The literature and repository comparison was targeted, not
exhaustive; independent assessment of historical novelty is still
needed. In particular, the article does not claim to settle a named
classical conjecture or all transseries differential equations.

\subsection{A compact statement of the main results}
Let \(\G\subset[0,\infty)\) be a nontrivial well-ordered additive
submonoid, with \(\delta=\min\G_{>0}>0\), and let \(A\) be a constant
complex matrix with real spectrum. Consider
\begin{equation}\label{eq:model-intro}
 (D-A)y=F(x,y),\qquad D=x\frac{d}{dx},\qquad
 F(x,z)=\sum_{\alpha,k}f_{\alpha,k}x^\alpha z^k,
\end{equation}
where \(\alpha\in\G\), \(k\in\NN^m\), and
\(f_{0,0}=f_{0,e_j}=0\). Set
\(\Reson=\spec(A)\cap\G_{>0}\).
For a polynomial vector \(P\), write \(\ell(P)=\max(0,\deg P)\),
with \(\ell(0)=0\). The normalized formal solution satisfies
\begin{equation}\label{eq:main-slope-intro}
 \Theta(y):=\sup_{\gamma\in\G_{>0}}
      \frac{\ell(P_\gamma)}{\gamma}
 =\max\left(\{0\}\cup
      \left\{\frac{\ell(P_\lambda)}{\lambda}:\lambda\in\Reson\right\}\right).
\end{equation}
Only actions in the finite set
\begin{equation}\label{eq:certificate-intro}
 \mathcal D=
 \bigcup_{\lambda\in\Reson}
 \{\alpha\in\G:\lambda-\alpha\in\G\}
\end{equation}
are needed to calculate the right-hand side. The exact slope is unchanged
by any logarithm-free local transformation with invertible constant
linear part.

For convergence, the additional universal condition is
\begin{equation}\label{eq:gap-intro}
 g(\G,A):=
 \inf_{\substack{\gamma\in\G_{>0},\ \lambda\in\spec(A)\\
                  \gamma\ne\lambda}}
 \abs{\gamma-\lambda}>0.
\end{equation}
When it holds, every absolutely convergent \(F\) has an absolutely
convergent normalized solution for every allowed collection of resonant
constants. When it fails, even logarithm-free linear forcing supplies
arbitrarily small counterexamples. The quantifier ``every'' is essential:
a particular nonlinear equation can have a convergent solution despite
\(g(\G,A)=0\).

\section{The Hahn--Dulac algebra and finite divisors}
\label{sec:support}

\subsection{Orientation, support, and polynomial blocks}
The independent variable tends to zero, so smaller positive exponents
are dominant. All well-orderings in this article use the usual order on
\(\RR\). This is opposite to the descending-monomial convention commonly
used at positive infinity.

Fix the monoid \(\G\) from the introduction. We use formal sums
\begin{equation}\label{eq:log-hahn}
 \mathscr D_\G=
 \left\{\sum_{\gamma\in\G}x^\gamma P_\gamma(L):
                  P_\gamma\in\CC[L]\right\},
 \qquad
 \mathfrak m_\G=
 \left\{y\in\mathscr D_\G:P_0=0\right\}.
\end{equation}
The coefficients at individual actions are finite polynomials, but their
degrees need not have a common bound. Thus \(\mathscr D_\G\) is larger
than the polynomial ring in \(L\) over an ordinary Hahn coefficient
ring. In the monomial ordering which first compares \(x\)-actions and
then logarithmic powers, its support remains well based: the first
nonempty action has only finitely many logarithmic terms.

Define a strong derivation by
\begin{equation}\label{eq:derivation}
 D(x^\gamma P(L))=x^\gamma\bigl(\gamma P(L)+P'(L)\bigr),
 \qquad DL=1.
\end{equation}
In the analytic sections, \(L=\Log x\) on a specified logarithmic cover.
The outer valuation \(\val(y)\) is the least action with nonzero
polynomial coefficient, and \(\val(0)=+\infty\). For vectors, use the
least action occurring in any component. This valuation ignores
logarithmic degree.

\begin{lemma}[Finite decompositions and substitution]\label{lem:finite}
For each \(\beta\in\G\), the set
\[
 \{(\alpha_1,\ldots,\alpha_n)\in\G^n:
               \alpha_1+\cdots+\alpha_n=\beta\}
\]
is finite for every fixed \(n\). If all entries are positive, then
\(n\le\floor{\beta/\delta}\).
Consequently, products in \(\mathscr D_\G\) and substitutions
\(F(x,y)\) with \(y\in\mathfrak m_\G^m\) are coefficientwise well defined.
\end{lemma}
\begin{proof}
For two summands this is Neumann's finite-decomposition lemma. An
elementary real-order proof is also useful: infinitely many distinct
pairs with fixed sum would have a subsequence whose first coordinates
strictly increase, forcing the second coordinates to strictly decrease,
contrary to well-ordering. Inducting on \(n\) gives the general case.
If every summand is positive, their sum is at least \(n\delta\).

At action \(\beta\), a monomial \(x^\alpha y^k\) can contribute only
when \(|k|\le\floor{\beta/\delta}\). There are finitely many such
multi-indices; for each, all action decompositions are finite. This also
allows \(\alpha=0\). Each contribution is a finite product of
polynomials in \(L\), so the resulting block is a polynomial.
\end{proof}

The same argument verifies the Leibniz rule for \(D\): every coefficient
comparison involves a finite sum. No interchange of analytic limits is
involved at this stage.

\begin{definition}[Divisor set]\label{def:divisor}
For \(\beta\in\G\), define
\begin{equation}\label{eq:divisor}
 \Div_\G(\beta)=\{\alpha\in\G:\beta-\alpha\in\G\}.
\end{equation}
An action \(\alpha\) in this set is called a divisor action of \(\beta\).
This is additive divisibility, not divisibility of real numbers.
\end{definition}

\begin{lemma}[Finite ancestry]\label{lem:divisor}
The set \(\Div_\G(\beta)\) is finite. If \(\alpha\) divides \(\beta\)
and \(\eta\) divides \(\alpha\), then \(\eta\) divides \(\beta\).
Every partial sum in a decomposition of \(\beta\) into monoid elements
is a divisor action of \(\beta\).
\end{lemma}
\begin{proof}
The map \(\alpha\mapsto(\alpha,\beta-\alpha)\) is a bijection to the
finite two-factor antidiagonal from Lemma~\ref{lem:finite}. Transitivity
follows from
\(\beta-\eta=(\beta-\alpha)+(\alpha-\eta)\in\G\).
The complementary partial sum proves the final assertion.
\end{proof}

\begin{remark}[Finite divisibility is not local finiteness]\label{rem:not-local}
For \(E=\{2-1/n:n\ge1\}\), the generated monoid
\(\G=\langle E\rangle\) is well ordered, although
\(\G\cap[0,2]\) is infinite. To see the well-ordering, each fixed
finite sumset of \(E\) is well ordered by Neumann's support lemma, and
a bounded interval meets only finitely many possible word lengths,
because \(\min E=1\). Nevertheless
\(\Div_\G(2)=\{0,1,2\}\).
This example is the main reason to formulate certificates with divisor
sets rather than with all actions below a cutoff.
\end{remark}

\section{The normalized nonlinear formal solution}
\label{sec:formal}

\subsection{Equation and normalization data}
Fix an integer \(m\ge1\), and write \(\NN=\{0,1,2,\ldots\}\).
Write the generalized eigenspace decomposition as
\[
 \CC^m=\bigoplus_{\lambda\in\spec(A)}V_\lambda,
 \qquad A|_{V_\lambda}=\lambda I+N_\lambda,
 \qquad N_\lambda^{s_\lambda}=0,
\]
where \(s_\lambda\) is the largest Jordan block size for \(\lambda\).
Let \(\Pi_\lambda\) denote the corresponding projection. The input is
\begin{equation}\label{eq:model}
 (D-A)y=F(x,y),\qquad
 F(x,z)=\sum_{\alpha\in\G,\,k\in\NN^m}
            f_{\alpha,k}x^\alpha z^k,
 \qquad f_{0,0}=f_{0,e_j}=0.
\end{equation}
The notation \(z^k\) means \(z_1^{k_1}\cdots z_m^{k_m}\).
Every \(f_{\alpha,k}\) is a vector in \(\CC^m\).
No convergence hypothesis is imposed yet.

The vanishing conditions have a precise role: constant forcing has
positive action, and all linear terms of action zero have already been
absorbed into \(A\). Pure nonlinear terms at action zero remain allowed.
For each \(\lambda\in\Reson=\spec(A)\cap\G_{>0}\), choose
\(c_\lambda\in V_\lambda\) and impose
\begin{equation}\label{eq:normalization}
 \Pi_\lambda P_\lambda(0)=c_\lambda.
\end{equation}
These are constants of the polynomial block at a resonant action; they
are not values of the analytic solution at \(x=0\).

\subsection{The polynomial Euler solver}
\begin{lemma}[One generalized eigenspace]\label{lem:solver}
Let \(N^s=0\), let \(Q\) be a polynomial vector, and consider
\begin{equation}\label{eq:block}
 (aI-N+\dd)P=Q.
\end{equation}
If \(a\ne0\), put \(B=(aI-N)^{-1}\). The unique polynomial solution is
\begin{equation}\label{eq:nonres-solver}
 P=\sum_{j=0}^{\deg Q}(-1)^j B^{j+1}\dd^jQ,
\end{equation}
with the zero solution when \(Q=0\). A nonzero \(Q\) and its solution
have the same degree.

If \(a=0\), the unique solution with \(P(0)=c\) is
\begin{equation}\label{eq:res-solver}
 P=e^{NL}c+\sum_{j=0}^{s-1}N^j I_L^{j+1}Q,
 \qquad I_LQ(L)=\int_0^LQ(t)\,dt.
\end{equation}
In this case
\(\ell(P)\le\max(s-1,\ell(Q)+s)\).
\end{lemma}
\begin{proof}
When \(a\ne0\), the inverse of \(aI-N\) exists by nilpotence.
The operators \(B\) and \(\dd\) commute, and differentiation is
nilpotent on every bounded-degree polynomial space. Multiplying the
finite geometric sum in \eqref{eq:nonres-solver} by
\((aI-N)(I+B\dd)\) gives \(Q\). Its leading coefficient is the
invertible image under \(B\) of the leading coefficient of \(Q\), so
the degree is unchanged. A nonzero polynomial in the kernel would have
a highest coefficient annihilated by \(aI-N\), which is impossible.

At resonance, the homogeneous solution with constant term \(c\) is
\(e^{NL}c\), a finite polynomial. Differentiating the sum in
\eqref{eq:res-solver} and subtracting its product by \(N\) telescopes:
all terms cancel except \(Q\), because \(N^s=0\). Every integral in
that sum vanishes at zero. The difference of two normalized solutions
solves \(P'=NP\) with zero constant term, and its successive
coefficients are therefore zero. Integration raises degree by one and
at most \(s\) integrations occur, proving the bound.
\end{proof}

Apply this lemma at action \(\gamma\) with \(a=\gamma-\lambda\) on
each generalized eigenspace. Let \(S\) be the resulting right inverse
of \(D-A\), with zero resonant constants. It does not introduce new
actions. Define the homogeneous seed
\begin{equation}\label{eq:seed}
 H_c=\sum_{\lambda\in\Reson}x^\lambda e^{N_\lambda L}c_\lambda.
\end{equation}
The desired equation is equivalent to
\begin{equation}\label{eq:fixedpoint}
 y=H_c+S F(x,y).
\end{equation}

\begin{theorem}[Existence, uniqueness, and finite coefficient dependence]
\label{thm:formal}
For every input \eqref{eq:model} and every collection
\((c_\lambda)_{\lambda\in\Reson}\), there is a unique solution
\(y\in\mathfrak m_\G^m\) satisfying \eqref{eq:normalization}.
Every individual block \(P_\gamma\) is obtained by a finite calculation
from earlier blocks and finitely many input coefficients.
No separation condition on the nonresonant divisors is required.
\end{theorem}
\begin{proof}
At action \(\gamma\), every occurrence of an unknown block in
\([x^\gamma]F(x,y)\) has action strictly below \(\gamma\).
For a positive-action linear term this follows by subtracting that
positive action. For a nonlinear term, at least one other factor has
positive action. Constant forcing has no unknown block. Thus
well-founded recursion on \(\G_{>0}\), using
Lemma~\ref{lem:solver}, constructs a polynomial block at every action.
Lemma~\ref{lem:finite} makes each step finite, and the prescribed
resonant constants make every step unique.

Equivalently, if \(u,v\) have positive action, expansion of a difference
of products gives
\begin{equation}\label{eq:valuation-gain}
 \val\bigl(F(x,u)-F(x,v)\bigr)
      \ge\val(u-v)+\delta.
\end{equation}
Each linear coefficient has positive action; each nonlinear difference
has at least one additional positive-action factor. Since \(S\) does
not lower actions, the iteration
\(y^{(0)}=0\), \(y^{(n+1)}=H_c+SF(x,y^{(n)})\) has differences whose
valuations tend to infinity. Hence each fixed coefficient stabilizes.
The support is contained in the already well-ordered set \(\G\), so
the stabilized coefficients define a valid formal sum. The same
valuation gain proves uniqueness. Arbitrarily small nonzero numbers
\(\gamma-\lambda\) cause no formal difficulty: each individual
polynomial calculation divides by only finitely many fixed nonzero
numbers.
\end{proof}

\begin{remark}[Which solutions are being classified]
The theorem classifies positive-action solutions supported in the chosen
\(\G\). A positive eigenvalue outside \(\G\) does not supply a free
homogeneous action in this class. One may enlarge the monoid to include
additional positive eigenvalues when appropriate. No assertion is made
about solutions with action zero, negative actions, inverse logarithms,
or other transseries scales.
\end{remark}

\section{A global degree budget and the exact resonance slope}
\label{sec:slope}

\subsection{A coarse budget that closes under nonlinear substitution}
For \(\gamma\in\G_{>0}\), define
\begin{equation}\label{eq:budget}
 b(\gamma)=\sum_{\substack{\lambda\in\Reson\\\lambda\le\gamma}}
                  s_\lambda,
 \qquad h(\gamma)=\floor{\gamma/\delta}\,b(\gamma),
 \qquad C=\frac{1}{\delta}\sum_{\lambda\in\Reson}s_\lambda.
\end{equation}
Set \(h(0)=0\). Empty sums are zero. These bounds depend only on the
ambient monoid and the constant matrix, not on the coefficients or the
chosen resonant constants.

\begin{lemma}[Superadditivity and a resonance margin]\label{lem:budget}
The function \(h\) is superadditive on \(\G\). If
\(\alpha+\gamma_1+\cdots+\gamma_n=\lambda\in\Reson\), where all
\(\gamma_i>0\), and the term is allowed by \eqref{eq:model}, then
every \(\gamma_i<\lambda\) and
\begin{equation}\label{eq:res-margin}
 \sum_i h(\gamma_i)\le h(\lambda)-s_\lambda.
\end{equation}
For \(n=0\), the same degree margin holds for constant forcing.
\end{lemma}
\begin{proof}
For positive \(u,v\), monotonicity of \(b\) and superadditivity of the
floor give
\[
 h(u)+h(v)
 \le\bigl(\floor{u/\delta}+\floor{v/\delta}\bigr)b(u+v)
 \le h(u+v).
\]
The case involving zero is immediate. At a resonant output action
\(\lambda\), the strict inequality \(\gamma_i<\lambda\) was proved
in Theorem~\ref{thm:formal}. Consequently
\(b(\gamma_i)\le b(\lambda)-s_\lambda\), and
\[
 \sum_i h(\gamma_i)
 \le\floor{\lambda/\delta}\bigl(b(\lambda)-s_\lambda\bigr)
 \le h(\lambda)-s_\lambda,
\]
because \(\floor{\lambda/\delta}\ge1\). Constant forcing has degree
zero, while the last expression is nonnegative.
\end{proof}

\begin{proposition}[Uniform linear degree bound]\label{prop:coarse}
The normalized solution satisfies
\begin{equation}\label{eq:coarse}
 \ell(P_\gamma)\le h(\gamma)\le C\gamma
 \qquad(\gamma\in\G_{>0}).
\end{equation}
In particular, no resonant action implies a logarithm-free solution.
\end{proposition}
\begin{proof}
Proceed by well-founded induction on the action. At a nonresonant
action, the degrees of all products in the source are bounded by the
sum of the input degree budgets, hence by \(h(\gamma)\); the
polynomial solver does not increase degree. At action
\(\lambda\in\Reson\), Lemma~\ref{lem:budget} bounds the source
degree by \(h(\lambda)-s_\lambda\). The resonant solver increases it
by at most \(s_\lambda\), and the homogeneous seed has degree at most
\(s_\lambda-1\le h(\lambda)\). Other generalized eigenspaces at the
same action use nonresonant solvers. This completes the induction.
The inequality \(h(\gamma)\le C\gamma\) follows directly from
\eqref{eq:budget}.
\end{proof}

The special resonance margin in \eqref{eq:res-margin} is important in
the analytic proof. A resonant inverse does \emph{not} map every
polynomial of degree at most \(h(\lambda)\) back into that same degree
space. It maps the smaller, naturally occurring source space into it.

\subsection{The exact finite-resonance maximum}
\begin{definition}[Maximal logarithmic slope]\label{def:slope}
For a positive-action vector series with finite block polynomials, put
\begin{equation}\label{eq:slope-definition}
 \Theta(y)=\sup_{\gamma\in\G_{>0}}
                     \frac{\ell(P_\gamma)}{\gamma}\in[0,\infty].
\end{equation}
Zero blocks contribute zero. This is a maximum over all actions when
it is attained; it is not a limit superior as the action tends to
infinity.
\end{definition}

\begin{theorem}[Exact resonance control]\label{thm:slope}
For the normalized solution of \eqref{eq:model},
\begin{equation}\label{eq:exact-slope}
 \boxed{\displaystyle
 \Theta(y)=\max\left(\{0\}\cup
  \left\{\frac{\ell(P_\lambda)}{\lambda}:\lambda\in\Reson\right\}\right).}
\end{equation}
Thus the optimal constant in
\(\ell(P_\gamma)\le\theta\gamma\) is attained at a resonant action
unless it is zero. The formula uses the actual resonant polynomials,
including all cancellations.
\end{theorem}
\begin{proof}
Let \(\theta\) denote the right-hand side, which is finite by
Proposition~\ref{prop:coarse}. We prove
\(\ell(P_\gamma)\le\theta\gamma\) by well-founded induction.
At a resonant action this is true by the definition of \(\theta\).
At a nonresonant action, each source product has degree at most
\[
 \sum_i\ell(P_{\gamma_i})
 \le\theta\sum_i\gamma_i
 \le\theta\gamma,
\]
since the coefficient action \(\alpha\) is nonnegative.
Logarithm-free forcing has degree zero. Sums cannot increase the
maximum degree, and the nonresonant solver preserves or lowers it.
This proves the upper bound at every action. The reverse inequality
follows by taking the supremum over the subset of resonant actions.
\end{proof}

\begin{corollary}[A finite test for logarithm-free solutions]
\label{cor:logfree}
The normalized solution is logarithm free if and only if
\(P_\lambda\) is constant for every \(\lambda\in\Reson\).
\end{corollary}
\begin{proof}
All degrees vanish if and only if \(\Theta(y)=0\), and
Theorem~\ref{thm:slope} characterizes that condition.
\end{proof}

\begin{remark}[What the maximum does not measure]
A finite solution \(y=xL\) has \(\Theta(y)=1\), even though it has no
large-action tail. Hence \eqref{eq:exact-slope} does not claim that
\(\deg P_\gamma/\gamma\) approaches the maximum infinitely often.
Nor does \(\Theta\) determine a convergence radius or coefficient
size. These distinctions are substantive, not just terminological.
\end{remark}

\section{Finite algebraic certificates and cancellation strata}
\label{sec:certificate}

Suppose \(\Reson\ne\varnothing\), let
\(\Lambda=\max\Reson\), and define \(\mathcal D\) by
\eqref{eq:certificate-intro}. When \(\Reson=\varnothing\), the slope
is identically zero and the certificate may be taken empty.

\begin{theorem}[Finite input certificate]\label{thm:certificate}
Fix \(A\) and \(\G\). All resonant polynomials \(P_\lambda\), and
therefore \(\Theta(y)\), depend only on the resonant constants and the
finite collection of source coordinates
\begin{equation}\label{eq:finite-input}
 f_{\alpha,k},\qquad
 \alpha\in\mathcal D,\quad |k|\le\floor{\Lambda/\delta}.
\end{equation}
Each coefficient of each \(P_\lambda\) is a polynomial in these
coordinates and the coordinates of \(c_\lambda\), with constants
depending on the fixed matrix and actions. The resonant polynomials can
be computed by recursion over the finite set \(\mathcal D\setminus\{0\}\).
\end{theorem}
\begin{proof}
The set \(\mathcal D\) is finite by Lemma~\ref{lem:divisor}. If a
source monomial contributes at action \(\beta\in\mathcal D\), each
of its input actions and its coefficient action divides \(\beta\).
Transitivity of divisibility places all of them in \(\mathcal D\).
Thus the coefficient recursion on \(\mathcal D\) never requires an
unknown block outside \(\mathcal D\). The arity bound in
Lemma~\ref{lem:finite} is at most \(\floor{\Lambda/\delta}\).

The source calculation uses only additions and multiplications. The
polynomial Euler solvers are linear in their source and normalization
constant, and divide only by fixed nonzero numbers such as
\(\gamma-\lambda\) and fixed positive integers. Induction over the
finite ordered set \(\mathcal D\setminus\{0\}\) proves polynomial
dependence. Including unused coordinates in \eqref{eq:finite-input}
does not change the assertion.
\end{proof}

\begin{theorem}[Algebraic slope strata]\label{thm:strata}
For fixed \(A\), \(\G\), and any real \(q\ge0\), the condition
\(\Theta(y)\le q\) is equivalent to the finite set of polynomial
equations
\begin{equation}\label{eq:strata}
 [L^d]P_\lambda=0
 \quad\text{for every }\lambda\in\Reson
 \text{ and integer }q\lambda<d\le h(\lambda).
\end{equation}
In particular, logarithm-free branches are characterized by finite
polynomial equations in the source parameters and resonant constants.
The possible slope values belong to the finite set
\begin{equation}\label{eq:possible-slopes}
 \{0\}\cup
 \left\{\frac d\lambda:
          \lambda\in\Reson,\ 1\le d\le h(\lambda),\ d\in\NN\right\}.
\end{equation}
\end{theorem}
\begin{proof}
By Theorem~\ref{thm:slope}, the inequality is exactly
\(\ell(P_\lambda)\le q\lambda\) at each resonance. The coarse
budget cuts off every polynomial at \(h(\lambda)\), so this is the
finite coefficient-vanishing condition \eqref{eq:strata}.
Theorem~\ref{thm:certificate} makes every displayed condition
polynomial in finite input coordinates. The exact maximum is either
zero or the ratio of an integer degree to a resonant action, which
gives \eqref{eq:possible-slopes}.
\end{proof}

The sublevel sets are algebraic in the finite parameter space. Exact
slope strata are obtained by subtracting smaller sublevel sets; no
claim that they are linear spaces is intended. An arbitrarily large
or infinite modification of source coordinates outside
\eqref{eq:finite-input} leaves the slope unchanged. Such a modification
can still change convergence, as Section~\ref{sec:accumulation} proves.

\subsection{A finite algorithm when divisor data are available}
A practical implementation receives an exact description of \(A\),
the resonances, and the finite divisor set \(\mathcal D\). Sort the
positive elements of \(\mathcal D\), compute the allowed source
coefficient at each action using already computed blocks, apply
\eqref{eq:nonres-solver} or \eqref{eq:res-solver}, and retain the
resonant blocks. Their coefficients give both the exact slope and the
polynomial equations for each stratum. This recursion terminates
because \(\mathcal D\) is finite, even if infinitely many ambient
actions lie below \(\Lambda\).

There is a necessary distinction between a \emph{finite mathematical
certificate} and a \emph{uniform procedure for discovering it}. With
arbitrary real exponents or an opaque infinite presentation of \(\G\),
we have not supplied an algorithm deciding membership, equality, or
additive divisibility. The finite routine is effective for exact
presentations that provide these data, such as a supplied rational
prefix known to be divisor closed. It must not silently replace that
hypothesis by an arbitrary finite truncation.

\section{Invariance under logarithm-free coordinate changes}
\label{sec:invariance}

Consider a transformation of dependent variables
\begin{equation}\label{eq:coordinate}
 y=H(x,z)=Bz+G(x,z),\qquad B\in\operatorname{GL}_m(\CC),
\end{equation}
where \(G\) has logarithm-free coefficients supported in \(\G\), and
its constant and linear terms of action zero vanish. Positive-action
translations and linear terms are allowed. No singular factor such as
\(x^a\) is permitted in the constant linear part.

\begin{lemma}[Formal inverse and slope monotonicity]\label{lem:coordinate}
The transformation \eqref{eq:coordinate} has a formal inverse of the
same type, with constant linear part \(B^{-1}\). For every
positive-action series \(z\) of finite slope,
\begin{equation}\label{eq:monotone-slope}
 \Theta(H(x,z))\le\Theta(z).
\end{equation}
\end{lemma}
\begin{proof}
For the inverse, solve
\(z=B^{-1}y-B^{-1}G(x,z)\) iteratively. Assign weight
\(\alpha+\delta|k|\) to a monomial \(x^\alpha y^k\).
Every term beyond the invertible constant linear part has positive
additional weight in differences: a positive coefficient action or
an additional dependent-variable factor contributes at least
\(\delta\). The same coefficient-stabilization argument as in
Theorem~\ref{thm:formal} therefore constructs the inverse. At a fixed
coefficient \((\alpha,k)\), finite decompositions and the weight
bound ensure that only finitely many terms occur. No logarithms are
introduced.

Let \(\theta=\Theta(z)\). A substituted product contributing at
action \(\gamma=\alpha+\sum_i\gamma_i\) has logarithmic degree at
most \(\theta\sum_i\gamma_i\le\theta\gamma\). The coefficient
itself has no logarithms. Sums and the constant matrix \(B\) do not
increase the bound, proving \eqref{eq:monotone-slope}.
\end{proof}

\begin{theorem}[Coordinate invariance of the exact slope]
\label{thm:invariant}
For every transformation \eqref{eq:coordinate} and every
positive-action series \(z\) of finite slope,
\begin{equation}\label{eq:invariance}
 \boxed{\Theta(H(x,z))=\Theta(z).}
\end{equation}
In particular, the finite-resonance maximum is a genuine invariant of
a normalized solution under this class of local coordinate changes.
\end{theorem}
\begin{proof}
Apply Lemma~\ref{lem:coordinate} first to \(H\), and then to its
logarithm-free inverse. The first application ensures the transformed
series still has finite slope. The second reverses the inequality.
\end{proof}

For compatible analytic transformations, this is also an invariant of
the convergent expansions of the transformed functions. Differentiating
\eqref{eq:coordinate} shows that the transformed differential equation
again has a logarithm-free nonlinear remainder and constant linear
matrix \(B^{-1}AB\). The theorem itself, however, only needs the formal
transformation and does not require convergence.

The restrictions are essential. A singular shearing
\(y=x^az\) shifts actions and can change every degree-to-action ratio.
A change of the independent variable rescales actions. A coordinate
change containing logarithms can insert them directly. None of those
operations belongs to \eqref{eq:coordinate}.

\section{A sharp universal criterion for analytic realization}
\label{sec:analytic}

\subsection{Absolute input convergence and the theorem}
Choose a Jordan basis once and use the vector \(\ell^1\)-norm in that
basis, with the induced matrix norm. Suppose that, for some \(R,S>0\),
\begin{equation}\label{eq:input-norm}
 \norm{F}_{R,S}:=
 \sum_{\alpha\in\G,\,k\in\NN^m}
       \norm{f_{\alpha,k}}R^\alpha S^{|k|}<\infty.
\end{equation}
This is an absolute coefficient condition, not merely a statement that
some function has the indicated formal expansion. It gives normally
convergent generalized powers in \(x\) and an analytic power series in
\(z\) on smaller domains.

For \(P(L)=\sum_dp_dL^d\), put
\begin{equation}\label{eq:polynorm}
 \norm{P}_T=\sum_d\norm{p_d}T^d,
 \qquad
 \norm{y}_{r,T}=\sum_{\gamma\in\G_{>0}}r^\gamma\norm{P_\gamma}_T,
 \qquad r,T>0.
\end{equation}
For vector-valued series the second norm is the sum of the component
norms. A solution is called \emph{coefficient-absolutely convergent}
if \(\norm{y}_{r,T}<\infty\) for some \(r,T>0\). We will prove that
this supplies actual convergent functions on sufficiently small
logarithmic sectors.

\begin{theorem}[Universal analytic realization]\label{thm:universal}
For fixed \(\G\) and \(A\), the following are equivalent.
\begin{enumerate}[label=\textup{(\roman*)}]
\item The nonresonant gap \(g(\G,A)\) in \eqref{eq:gap-intro} is positive.
\item For every \(F\) satisfying \eqref{eq:model} and
\eqref{eq:input-norm}, and every allowed collection of resonant
constants, the normalized formal solution is coefficient-absolutely
convergent.
\end{enumerate}
If the gap is positive, explicit sufficient norm inequalities for a
convergent solution are given in \eqref{eq:contraction-conditions}.
If it vanishes, counterexamples can be chosen with no dependence on
\(y\), with all resonant constants zero, and with arbitrarily small
input norm. Their formal solutions are logarithm free.
\end{theorem}

The universal quantifier is over all coefficient series supported in
\(\G\), not just those supported in a selected proper generating set.
The theorem concerns the specific normalized expansion; it does not
say that a differential equation has no other analytic interpretation
when that expansion diverges.

\subsection{The degree-budget Banach space}
For fixed \(r,T\), let \(\mathcal X_{r,T}\) be the space of series
with positive actions, finite norm \eqref{eq:polynorm}, and
\(\deg P_\gamma\le h(\gamma)\). Zero polynomials are always allowed.
It is a Banach space: it is a closed coordinate subspace of the
weighted \(\ell^1\)-space indexed by the pairs
\((\gamma,d)\), \(0\le d\le h(\gamma)\).
The scalar version is closed under multiplication, by
superadditivity of \(h\), and
\begin{equation}\label{eq:norm-product}
 \norm{uv}_{r,T}\le\norm{u}_{r,T}\norm{v}_{r,T}.
\end{equation}

We need a slightly smaller source space at resonances. Let
\(\mathcal Q_{r,T}\) consist of positive-action series \(Q\) with
\[
 \deg Q_\gamma\le h(\gamma)\quad(\gamma\notin\Reson),
 \qquad
 \deg Q_\lambda\le h(\lambda)-s_\lambda
                    \quad(\lambda\in\Reson).
\]
These inequalities apply to the full vector polynomial. The source
\(F(x,y)\) belongs to \(\mathcal Q_{r,T}\) whenever the substitution
converges and \(y\in\mathcal X_{r,T}\), by the resonance margin in
Lemma~\ref{lem:budget}. The same holds for a difference of two such
sources. This domain restriction prevents an illicit use of a bounded
resonant inverse on a degree space that it does not preserve.

\subsection{An explicit uniform Euler-inverse bound}
Define
\begin{equation}\label{eq:analytic-constants}
 \Lambda_0=\max_{\lambda\in\spec(A)}|\lambda|,
 \quad M=\max_{\lambda\in\spec(A)}\norm{N_\lambda},
 \quad B_0=4(\Lambda_0+M+1).
\end{equation}
Choose
\begin{equation}\label{eq:T-choice}
 T\ge\max(1,8C),\qquad D_0=\ceil{CB_0}.
\end{equation}
Assume \(g=g(\G,A)>0\), and set
\begin{align}
 M_0&=\max_{\lambda\in\spec(A)}
      \sum_{j=0}^{s_\lambda-1}
         \norm{N_\lambda}^{j}g^{-(j+1)},
       \label{eq:Mzero}\\
 K_{\mathrm{low}}&=
      \sum_{j=0}^{D_0}M_0^{j+1}
           \frac{D_0!}{(D_0-j)!\,T^j},
       \label{eq:Klow}\\
 K_{\mathrm{res}}&=
      \max\left(\{0\}\cup
      \left\{\sum_{j=0}^{s_\lambda-1}
           \norm{N_\lambda}^{j}\frac{T^{j+1}}{(j+1)!}
           :\lambda\in\Reson\right\}\right),
       \label{eq:Kres}\\
 K&=\max\{K_{\mathrm{low}},K_{\mathrm{res}},4/B_0\}.
       \label{eq:K}
\end{align}
Powers of a matrix norm with exponent zero are interpreted as one,
also when the matrix is zero. These are deliberately conservative
constants; they are useful because all dependence is explicit.

\begin{lemma}[Bounded right inverse]\label{lem:bounded}
For every \(r>0\), the zero-normalized polynomial solver defines a
bounded linear map
\[
 S:\mathcal Q_{r,T}\longrightarrow\mathcal X_{r,T},
 \qquad \norm{SQ}_{r,T}\le K\norm{Q}_{r,T}.
\]
\end{lemma}
\begin{proof}
All estimates are coefficientwise in the action. On polynomials of
degree at most \(d\),
\begin{equation}\label{eq:derivative-norm}
 \norm{\dd P}_T\le\frac dT\norm{P}_T,
 \qquad
 \norm{\dd^jP}_T\le
     \frac{d!}{(d-j)!\,T^j}\norm{P}_T
       \quad(0\le j\le d).
\end{equation}
The integrals with zero constant term satisfy
\begin{equation}\label{eq:integral-norm}
 \norm{I_L^kP}_T\le\frac{T^k}{k!}\norm{P}_T,
\end{equation}
since the coefficient multiplier is
\(d!/(d+k)!\le1/k!\).

For a nonresonant block with action \(\gamma<B_0\), the finite
nilpotent inverse has norm
\[
 \norm{((\gamma-\lambda)I-N_\lambda)^{-1}}
 \le\sum_{j=0}^{s_\lambda-1}
        \norm{N_\lambda}^{j}|\gamma-\lambda|^{-(j+1)}
 \le M_0.
\]
The degree is at most \(C\gamma\le D_0\). Inserting
\eqref{eq:derivative-norm} into \eqref{eq:nonres-solver} gives
\(K_{\mathrm{low}}\).

For \(\gamma\ge B_0\), put \(a=\gamma-\lambda\). We have
\(a\ge3\gamma/4\), \(\norm{N_\lambda}\le\gamma/4\), and
\(\norm{\dd}\le h(\gamma)/T\le\gamma/8\). Thus, on the finite
polynomial space,
\[
 \norm{(aI+\dd-N_\lambda)^{-1}}
 \le\frac{1}{a-\norm{\dd-N_\lambda}}
 \le\frac{8}{3\gamma}\le\frac4\gamma\le\frac4{B_0}.
\]
The displayed inverse follows from a norm-convergent Neumann series;
it equals the polynomial inverse because that inverse is unique.
No resonance occurs at these large actions, since all eigenvalues have
absolute value at most \(\Lambda_0<B_0\).

At a resonant block, \eqref{eq:res-solver} with zero constant and
\eqref{eq:integral-norm} give \(K_{\mathrm{res}}\). The source
margin makes its output degree at most \(h(\lambda)\).
Other eigenspaces at that action are covered by the nonresonant
estimate. In a direct-sum \(\ell^1\)-norm, the maximum of the block
bounds bounds the full operator. Finally multiply by \(r^\gamma\)
and sum over actions. Absolute summability justifies this last step.
\end{proof}

The large-action inverse compensates for the growth of polynomial
degree; the finite-action gap controls possible small divisors. These
are separate parts of the proof. The method is an adaptation of the
classical polynomial-norm approach acknowledged in
Section~\ref{sec:introduction}, rather than a claim that weighted
Dulac majorants are themselves new.

\subsection{Quantitative contraction and sufficiency}
For \(0<r<R\) and \(0<\eta<S\), define
\begin{align}
 \Phi_0(r)&=\sum_{\alpha>0}\norm{f_{\alpha,0}}r^\alpha,
                  \label{eq:forcing-majorant}\\
 L_F(r,\eta)&=
   \sum_{\substack{\alpha\in\G,\,k\in\NN^m\\|k|\ge1}}
   |k|\norm{f_{\alpha,k}}r^\alpha\eta^{|k|-1}.
                  \label{eq:lipschitz-majorant}
\end{align}
These quantities are finite on smaller radii. In the ball
\(\norm{y}_{r,T},\norm{z}_{r,T}\le\eta\), the usual telescoping
product identity and \eqref{eq:norm-product} imply
\begin{equation}\label{eq:majorant-Lipschitz}
 \norm{F(x,y)-F(x,z)}_{r,T}
 \le L_F(r,\eta)\norm{y-z}_{r,T}.
\end{equation}
The finite homogeneous seed satisfies
\begin{equation}\label{eq:seed-norm}
 \norm{H_c}_{r,T}
 \le\sum_{\lambda\in\Reson}
        r^\lambda e^{\norm{N_\lambda}T}\norm{c_\lambda}.
\end{equation}

\begin{proposition}[A usable contraction certificate]
\label{prop:contraction}
Choose \(T,K\) as above. Suppose \(0<r<R\), \(0<\eta<S\), and
\(0\le\kappa<1\) satisfy
\begin{equation}\label{eq:contraction-conditions}
 \boxed{\begin{aligned}
 K L_F(r,\eta)&\le\kappa,\\
 \norm{H_c}_{r,T}+K\Phi_0(r)&\le(1-\kappa)\eta.
 \end{aligned}}
\end{equation}
Then the normalized solution belongs to \(\mathcal X_{r,T}\) and
has norm at most \(\eta\). Such choices exist for every input
\eqref{eq:input-norm} and every fixed collection \(c_\lambda\).
\end{proposition}
\begin{proof}
On the closed radius-\(\eta\) ball in \(\mathcal X_{r,T}\), define
\(\mathcal T(y)=H_c+SF(x,y)\). The source lies in
\(\mathcal Q_{r,T}\), as proved above. Lemma~\ref{lem:bounded} and
\eqref{eq:majorant-Lipschitz} show that \(\mathcal T\) has
Lipschitz constant at most \(\kappa\). Moreover
\[
 \norm{\mathcal T(y)}_{r,T}
 \le\norm{H_c}_{r,T}+K\Phi_0(r)
              +K L_F(r,\eta)\eta\le\eta.
\]
The Banach contraction theorem gives a unique fixed point in this
ball. Its coefficient identities and resonant constants are exactly
those of Theorem~\ref{thm:formal}, so it is the normalized formal
solution already constructed.

To choose the radii, first make \(\eta>0\) small. The action-zero
part of \(L_F(r,\eta)\) then tends to zero because it begins at
nonlinear degree two. For this fixed \(\eta<S\), every
positive-action part of \(L_F(r,\eta)\) tends to zero as
\(r\downarrow0\), by dominated convergence using the absolute input
norm. The same holds for \(\Phi_0(r)\) and the finite positive-action
seed norm. Choose \(r\) small enough to satisfy both inequalities,
for example with \(\kappa=1/2\). Derivative majorants converge at
\(\eta<S\) because \(n(\eta/S)^{n-1}\) is bounded. This verifies
all uses of dominated convergence and completes the existence proof.
\end{proof}

This proves the sufficient direction of
Theorem~\ref{thm:universal}. Notice that the resonant constants need
not be small: their positive actions make the seed small after the
radius in \(x\) is reduced.

\subsection{Necessity and the direction of accumulation}
\begin{lemma}[Gap failure is left accumulation]\label{lem:gap-failure}
The gap \(g(\G,A)\) is zero if and only if some positive eigenvalue
\(\lambda\) is a left accumulation point of \(\G\).
The eigenvalue itself need not belong to \(\G\).
\end{lemma}
\begin{proof}
Left accumulation gives arbitrarily small nonzero distances and hence
a zero gap. Conversely, a sequence of distances tending to zero has,
after passage to a subsequence, one fixed eigenvalue because the
spectrum is finite. Nonpositive eigenvalues cannot occur: positive
actions are bounded below by \(\delta\). Infinitely many distinct
actions approaching the eigenvalue from above would contain a
strictly decreasing subsequence, impossible in a well-ordered set.
The remaining subsequence approaches from below and may be chosen
strictly increasing.
\end{proof}

\begin{proof}[Proof of necessity in Theorem~\ref{thm:universal}]
Let \(\lambda>0\) be a left accumulation point as in the lemma.
Choose distinct \(\gamma_n\uparrow\lambda\), with
\(\varepsilon_n=\lambda-\gamma_n\le2^{-n}\), and an eigenvector
\(v\ne0\) satisfying \(Av=\lambda v\). For any nonzero scalar
\(\tau\), take the source
\begin{equation}\label{eq:gap-counterexample}
 F(x,z)=\tau\sum_{n\ge1}\frac{\varepsilon_n}{n}
                       x^{\gamma_n}v.
\end{equation}
It is independent of \(z\) and has finite absolute input norm at
every fixed \(R,S>0\): the actions stay in a bounded positive
interval, while \(\sum\varepsilon_n/n\) converges. Its norm can be
made arbitrarily small by reducing \(|\tau|\).

With all resonant constants zero, the coefficient at each
\(\gamma_n\) is the constant vector \(-\tau v/n\), since
\[
 (\gamma_n I-A)(-\tau v/n)
       =\tau\varepsilon_n v/n.
\]
Thus \(\Theta(y)=0\), but for every \(r>0\),
\[
 \norm{y}_{r,T}
 \ge |\tau|\norm v\sum_{n\ge1}\frac{r^{\gamma_n}}n=\infty,
\]
because \(r^{\gamma_n}\to r^\lambda>0\). The same subseries
prevents pointwise absolute convergence at every positive \(x\).
Therefore the universal property fails. This proves necessity and
all the counterexample assertions.
\end{proof}

\begin{corollary}[Two independent obstructions]\label{cor:independent}
Finite resonance data determines the exact logarithmic slope without
any gap assumption. A zero logarithmic slope does not imply analytic
realization. Conversely, a vanishing universal gap does not force
every particular input to diverge.
\end{corollary}
\begin{proof}
The first two assertions follow from Theorem~\ref{thm:slope} and
\eqref{eq:gap-counterexample}. For the last assertion, the zero source
with zero constants gives a convergent zero solution even when the
gap vanishes. Section~\ref{sec:accumulation} gives a much less
degenerate nonlinear family with an exact convergence criterion.
\end{proof}

\section{Actual functions and explicit remainder bounds}
\label{sec:evaluation}

Let \(y\) have finite coefficient norm \(\norm y_{r,T}=M_y\), and
suppose its degree bound is \(\ell(P_\gamma)\le\theta\gamma\).
For a normalized solution one may use the exact
\(\theta=\Theta(y)\), although the coarse constant \(C\) is available
without computing the certificate. On a chosen logarithmic cover put
\begin{equation}\label{eq:q}
 q_\theta(x)=\frac{|x|}{r}
        \max\left(1,\frac{|\Log x|}{T}\right)^\theta.
\end{equation}

\begin{theorem}[Evaluation and action-tail certificate]
\label{thm:remainder}
On the domain \(q_\theta(x)<1\), the series
\(\sum_\gamma x^\gamma P_\gamma(\Log x)\) converges absolutely and
normally on compact subsets. For every \(B>0\),
\begin{equation}\label{eq:tail}
 \boxed{\displaystyle
 \norm{\sum_{\gamma\ge B}x^\gamma P_\gamma(\Log x)}
       \le M_y q_\theta(x)^B.}
\end{equation}
The series may be differentiated term by term with \(D=x\,d/dx\).
When it comes from Proposition~\ref{prop:contraction}, it solves the
actual analytic equation and one may replace \(M_y\) by \(\eta\).
Every sector with bounded argument contains a sufficiently small
punctured subsector on which these conclusions hold.
\end{theorem}
\begin{proof}
For a monomial with \(d\le\theta\gamma\),
\[
 |x|^\gamma|\Log x|^d
 \le r^\gamma T^d
   \left(\frac{|x|}{r}
       \max(1,|\Log x|/T)^\theta\right)^\gamma.
\]
Multiply by \(\norm{p_{\gamma,d}}\), sum over \(d\), and then sum
over actions. Since \(q^\gamma\le1\), absolute convergence follows
from the coefficient norm. Restricting to \(\gamma\ge B\) gives
\(q^\gamma\le q^B\), which proves \eqref{eq:tail}.
On a compact subset of \(q_\theta<1\), the maximum of \(q_\theta\)
is some \(q_*<1\), so the same majorant proves normal convergence.

For differentiation, the action part contributes a factor \(\gamma\).
The logarithmic derivative contributes at most \(d/T\) to the same
weighted monomial majorant, because
\[
 \frac{|\Log x|^{d-1}}{T^d}
 \le\frac1T\max(1,|\Log x|/T)^d
 \quad(d\ge1).
\]
Thus the differentiated coefficient is bounded by
\(\gamma(1+\theta/T)q_*^\gamma\) times its coefficient-norm weight.
The function \(\gamma q_*^\gamma\) is bounded on
\([\delta,\infty)\), proving local normal convergence of the
derivative. Analytic substitution in \(F\) is permitted by its
absolute majorant and \(\norm{y(x)}\le M_y\le\eta<S\).
Coefficient identities therefore become the analytic differential
identity.

Finally, on a bounded-argument sector,
\(|x|\max(1,|\Log x|/T)^\theta\to0\) as \(x\to0\).
This proves the sector assertion.
\end{proof}

\begin{corollary}[Finite asymptotic prefixes]\label{cor:prefix}
For a convergent solution as above, any finite initial collection of
nonzero action blocks is an asymptotic prefix of the actual function.
More precisely, if the next action after that collection is \(\beta\),
the remaining sum is
\(O(|x|^\beta\max(1,|\Log x|)^ {\theta\beta})\)
on a bounded-argument sector. It is therefore smaller than every
nonzero logarithmic monomial in the last retained action block.
\end{corollary}
\begin{proof}
Use \eqref{eq:tail} with \(B=\beta\). The next action is strictly
larger than the last retained one, and a positive power of \(|x|\)
dominates any fixed power of \(|\Log x|\) as \(x\to0\).
\end{proof}

\begin{remark}[A cutoff need not be a finite computation]
The sum over \(\gamma<B\) may itself contain infinitely many blocks.
Theorem~\ref{thm:remainder} is a bound on an entire action tail, not a
promise that ordinary finite truncation has reached that tail. In the
accumulating example, reaching action two requires either summing or
bounding an infinite block of actions below two. This distinction must
be retained in numerical implementations.
\end{remark}

\subsection{A posteriori certification inside the norm ball}
The contraction also provides a residual certificate that is useful
when coefficients are computed only approximately or whole action
blocks are represented by separate analytic routines.

\begin{proposition}[Fixed-point residual bound]\label{prop:residual}
Under \eqref{eq:contraction-conditions}, let
\(z\in\mathcal X_{r,T}\) have norm at most \(\eta\), and suppose
\[
 e\ge\norm{z-H_c-SF(x,z)}_{r,T}.
\]
Then
\begin{equation}\label{eq:residual}
 \norm{y-z}_{r,T}\le\frac{e}{1-\kappa}.
\end{equation}
For pointwise evaluation, the coarse degree bound gives
\[
 \norm{y(x)-z(x)}
 \le q_C(x)^\delta\frac{e}{1-\kappa}
 \qquad(q_C(x)<1).
\]
\end{proposition}
\begin{proof}
Write \(\mathcal T=H_c+SF\). Since \(y=\mathcal T(y)\),
\[
 \norm{y-z}\le\norm{\mathcal T(y)-\mathcal T(z)}
                   +\norm{\mathcal T(z)-z}
 \le\kappa\norm{y-z}+e.
\]
Rearranging proves \eqref{eq:residual}. The difference has positive
actions and degree at most \(h(\gamma)\le C\gamma\). The evaluation
estimate with its smallest possible action \(\delta\) proves the
pointwise bound.
\end{proof}

The word ``certificate'' here means a theorem with explicitly stated
inequalities. A numerical application must actually bound the source
majorants and residual; unverified floating-point estimates are not
a substitute for those hypotheses.

\section{Explicit nonlinear examples}
\label{sec:examples}

\subsection{An unbounded logarithmic degree with an exact invariant cone}
Consider the scalar equation
\begin{equation}\label{eq:geometric-ode}
 (D-1)y=y^2+x(1+2y+y^2),\qquad \G=\NN.
\end{equation}
The logarithm-free change \(z=y/(1+y)\) has an invertible constant
linear part. Direct differentiation gives
\[
 Dz=\frac{Dy}{(1+y)^2}=z+x.
\]
Thus
\begin{equation}\label{eq:geometric-solution}
 z=x(L+c),\qquad
 y=\frac{x(L+c)}{1-x(L+c)}
   =\sum_{n\ge1}x^n(L+c)^n.
\end{equation}
The resonant action is one, its normalized constant is \(c\), and
\(\deg P_n=n\) for every \(n\ge1\). Therefore
\(\Theta(y)=1=\Theta(z)\). The geometric series converges exactly
when \(|x(\Log x+c)|<1\).

This example simultaneously proves three limitations of a naive
extension of linear theory. Nonlinear solutions need not have a
uniform bound on logarithmic degree. They need not belong to
\(\CC((x^\G))[L]\), although they use only one logarithm. And a
uniform finite degree bound is not preserved by invertible
logarithm-free coordinate changes: \(z\) has degree one, while \(y\)
has unbounded degrees. The slope, in contrast, is preserved.

\subsection{An algebraic locus of logarithm-free branches}
Consider
\begin{equation}\label{eq:cancellation-example}
 Du=u,\qquad (D-2)v=u^2+a x^2.
\end{equation}
The positive-action normalized solutions are
\begin{equation}\label{eq:cancellation-solution}
 u=cx,\qquad v=x^2\bigl(d+(c^2+a)L\bigr).
\end{equation}
Consequently
\[
 \Theta(u,v)=
 \begin{cases}
 0,&c^2+a=0,\\
 1/2,&c^2+a\ne0.
 \end{cases}
\]
The logarithm-free condition is a genuine nonlinear algebraic
constraint on a homogeneous parameter, not a vector-space kernel.
Over the reals, \(a>0\) admits no logarithm-free branch, \(a<0\)
admits two values of \(c\), and \(a=0\) forces \(c=0\).
The parameter \(d\) remains free. Over \(\CC\), the same polynomial
equation describes the locus without an order restriction.

\subsection{A resonant cascade}
For \(r\ge1\), consider
\begin{equation}\label{eq:cascade}
 (D-1)y_1=x,\qquad
 (D-j)y_j=y_1y_{j-1}\quad(2\le j\le r),
\end{equation}
with all resonant constants zero. Then
\begin{equation}\label{eq:cascade-solution}
 y_j=\frac{x^jL^{2j-1}}{(2j-1)!!}.
\end{equation}
For \(j=1\), this is \(y_1=xL\). If the formula holds for
\(j-1\), the source at action \(j\) is
\(L^{2j-2}/(2j-3)!!\); a single integration in \(L\) gives
\eqref{eq:cascade-solution}. Hence
\begin{equation}\label{eq:cascade-slope}
 \Theta(y_1,\ldots,y_r)=\max_{1\le j\le r}\frac{2j-1}{j}
                    =2-\frac1r.
\end{equation}
This is a finite cascade of nonlinear resonant integrations.
The largest slope is detected at the final resonant action; the
calculation does not require estimating an infinite tail.

\subsection{A Jordan block feeding a nonlinear resonance}
Let
\begin{equation}\label{eq:jordan-example}
 Du=u+v,\qquad Dv=v+x,\qquad Dw=2w+u^2,
\end{equation}
again with zero resonant constants. The solution is
\begin{equation}\label{eq:jordan-solution}
 v=xL,\qquad u=\frac{xL^2}{2},\qquad
 w=\frac{x^2L^5}{20}.
\end{equation}
Indeed, the two-dimensional Jordan block integrates the forcing twice
across its components, and the last equation integrates
\(L^4/4\) at action two. The exact slope is \(5/2\).
The example illustrates why the coarse budget must account for Jordan
block sizes, even though the exact theorem only reads off the
resulting resonant polynomial degrees.

\section{An exact nonlinear accumulation threshold}
\label{sec:accumulation}

The universal theorem treats all sources on a fixed monoid. A sharper
input-specific statement is possible when the accumulation lies below
the first nonlinear interaction. It gives a direct nonlinear extension
of the linear \(p>2\) phenomenon recorded in the repository's
Hahn--Fuchsian article \cite{repo-fuchsian}.

\subsection{The monoid and the equation}
Let
\begin{equation}\label{eq:accum-monoid}
 E=\{2-1/n:n\ge1\},\qquad \G=\langle E\rangle,
 \qquad \delta=1.
\end{equation}
Below action two, the only positive actions are those in \(E\).
Every sum of two positive actions is at least two, and equality holds
only for \(1+1\). Moreover
\begin{equation}\label{eq:accum-divisors}
 \Div_\G(2)=\{0,1,2\},\qquad
 \min\bigl(\G\cap(2,\infty)\bigr)=\frac52.
\end{equation}
The final identity follows because the smallest positive generator
after one is \(3/2\): a sum of two generators other than \(1+1\)
is at least \(5/2\), and sums of three or more are at least three.

Consider the scalar nonlinear equation
\begin{equation}\label{eq:accum-ode}
 (D-2)y=a x+b x^2+c y^2+
          \sum_{n=2}^{\infty}t_n x^{2-1/n},
 \qquad \sum_{n=2}^{\infty}|t_n|<\infty,
\end{equation}
with normalization \(P_2(0)=k\). All four constants
\(a,b,c,k\) are arbitrary complex numbers. The input satisfies the
absolute coefficient condition at every fixed finite \(x\)-radius
and every finite \(y\)-radius, because it is quadratic in \(y\)
and the actions of the tail stay in a bounded interval.

\begin{theorem}[Finite slope and infinite summability are independent]
\label{thm:sharp-family}
The normalized formal solution of \eqref{eq:accum-ode} has
\begin{align}
 P_1&=-a,\label{eq:family-P1}\\
 P_{2-1/n}&=-nt_n\qquad(n\ge2),\label{eq:family-Pn}\\
 P_2&=k+(b+ca^2)L.\label{eq:family-P2}
\end{align}
Its exact slope is
\begin{equation}\label{eq:family-slope}
 \boxed{\displaystyle
 \Theta(y)=
 \begin{cases}
 0,&b+ca^2=0,\\
 1/2,&b+ca^2\ne0.
 \end{cases}}
\end{equation}
Independently of these constants, the normalized solution is
coefficient-absolutely convergent if and only if
\begin{equation}\label{eq:family-threshold}
 \boxed{\displaystyle \sum_{n=2}^{\infty}n|t_n|<\infty.}
\end{equation}
Thus the entire accumulation tail is invisible to the exact slope,
but it decides convergence through a sharp weighted summability test.
\end{theorem}

\begin{proof}[Proof of the coefficient and slope assertions]
There is no quadratic interaction at action below two. At action one,
\((1-2)P_1=a\), giving \eqref{eq:family-P1}. At action \(2-1/n\),
\((-1/n+\dd)P_{2-1/n}=t_n\), whose unique polynomial solution is
\eqref{eq:family-Pn}. At action two, the only product contributing
from \(y^2\) is \(P_1^2=a^2\); consequently
\(\dd P_2=b+ca^2\). The normalization gives
\eqref{eq:family-P2}. Since two is the only resonant action,
Theorem~\ref{thm:slope} gives \eqref{eq:family-slope}.
The divisor calculation \eqref{eq:accum-divisors} explains the finite
certificate directly: none of the actions \(2-1/n\), \(n\ge2\),
can be an ancestor of the action-two resonance.
\end{proof}

\begin{proof}[Proof of necessity of the summability condition]
For every \(r>0\), the numbers \(r^{2-1/n}\) tend to \(r^2>0\).
The constant coefficients in \eqref{eq:family-Pn} therefore have
finite weighted norm only when \(\sum n|t_n|<\infty\).
The same reasoning proves failure of pointwise absolute convergence
at every positive \(x\) when this sum diverges. Adding a resonant
homogeneous term cannot cancel these coefficients at distinct
nonresonant actions.
\end{proof}

\begin{proof}[Proof of sufficiency, including nonlinear feedback]
Assume \(\sum n|t_n|<\infty\), put \(q=b+ca^2\), and form the
entire low-action block
\begin{equation}\label{eq:low-core}
 Y_0=-a x-\sum_{n=2}^{\infty}nt_nx^{2-1/n}
                      +x^2(k+qL).
\end{equation}
This is an infinite, but coefficient-absolutely convergent, core.
For any fixed \(T>0\), its norm tends to zero as \(r\downarrow0\).
Its degree slope is \(\theta=0\) when \(q=0\), and \(\theta=1/2\)
otherwise. Write \(y=Y_0+w\), where \(w\) has actions greater than
two. Direct computation using \eqref{eq:low-core} reduces the equation
to
\begin{equation}\label{eq:tail-equation}
 (D-2)w=c\bigl((Y_0+w)^2-a^2x^2\bigr).
\end{equation}
The right side has no actions at most two: the coefficient of
\(Y_0^2\) at two is exactly \(a^2\), and every other product has
action at least \(5/2\).

Use the coefficient-norm space of series supported in
\(\G\cap(2,\infty)\) with
\(\deg P_\gamma\le\floor{\theta\gamma}\). It is complete, and the
right side of \eqref{eq:tail-equation} respects this degree bound.
On these actions, \(\gamma-2\ge1/2\). The nonresonant part of the
proof of Lemma~\ref{lem:bounded}, with linear degree constant
\(\theta\), therefore gives a bounded inverse \(S_{>}\) of
\(D-2\) on this tail space. Let its norm bound be \(K_{>}\).
For a ball \(\norm w\le\eta\), the map
\[
 \mathcal T_{>}(w)=c S_{>}
                   \bigl((Y_0+w)^2-a^2x^2\bigr)
\]
has Lipschitz constant at most
\(2|c|K_{>}(\norm{Y_0}+\eta)\), and its value at zero has norm at
most
\( |c|K_{>}(\norm{Y_0}^2+|a|^2r^2)\).
First choose \(\eta>0\) small enough, and then \(r>0\) small enough,
to make the Lipschitz constant less than \(1/2\) and the value at
zero at most \(\eta/2\). When \(c=0\), simply take \(w=0\).
The contraction theorem gives a convergent tail \(w\).

The resulting \(Y_0+w\) has the prescribed resonant constant and
satisfies the formal equation. By Theorem~\ref{thm:formal}, it is the
normalized formal solution. This proves sufficiency without assuming
a gap below action two; all small divisors there have been absorbed
into the summable infinite core.
\end{proof}

\begin{corollary}[The exact power-law threshold]\label{cor:p-threshold}
For \(t_n=\varepsilon n^{-p}\), with
\(\varepsilon\ne0\) and real \(p>1\), the input of
\eqref{eq:accum-ode} is coefficient-absolutely convergent and its
normalized solution is coefficient-absolutely convergent exactly when
\(p>2\). This threshold is independent of \(a,b,c,k\).
\end{corollary}
\begin{proof}
The input requires \(\sum n^{-p}<\infty\). Theorem~\ref{thm:sharp-family}
requires \(\sum n^{1-p}<\infty\) for the solution. These are the
stated strict inequalities.
\end{proof}

\subsection{Consequences and limitations of the example}
Fix \(a,b,c,k\), and compare the tails
\(t_n=\varepsilon n^{-2}\) and \(t_n=\varepsilon n^{-3}\).
They agree in every coordinate of the finite resonance certificate,
because that certificate does not contain any \(t_n\).
Their slopes are identical, yet the first normalized expansion
diverges and the second converges. Both tails can be arbitrarily
small in the input norm. Taking \(b=-ca^2\) makes both formal
solutions logarithm free, so absence of logarithms does not weaken
the distinction.

The convergence assertion here is stronger than merely applying the
universal theorem: the universal gap is zero for this monoid and
matrix. The proof instead removes a summable infinite low-action
block and uses the positive gap in the remaining tail. This is a
precise instance of an infinite-core strategy, not a claim that every
accumulating equation admits such a reduction.

When the weighted sum diverges, the theorem excludes the stated
coefficient expansion as an absolutely convergent representation. It
does not exclude solutions defined by another normalization, a
renormalized summation procedure, or direct integration away from
zero. Establishing canonical replacements in nonlinear cases is a
separate research problem.

\section{Computation, reproducibility, and a formalization route}
\label{sec:verification}

\subsection{What the accompanying program actually checks}
The package contains \texttt{verification/verify.py} and its recorded
\texttt{results.json}. The program uses exact SymPy polynomial and
rational arithmetic for the finite algebraic checks. Its scalar and
Jordan solver uses descending component substitution, providing a
second implementation of the operator identities in
Lemma~\ref{lem:solver}. A finite action dictionary implements
coefficient convolution and the normalized nonlinear recursion.

The recorded run used Python 3.13.5, SymPy 1.14.0, mpmath 1.3.0, and
random seed 20260929. All checks passed. The categories below are
reported separately rather than added into a purported count of
independent theorem verifications.

\begin{center}
\begin{tabular}{@{}p{0.72\linewidth}r@{}}
\toprule
Check category & Recorded count\\
\midrule
Exact scalar/Jordan block component identities & 200\\
Exact nonlinear system coefficient-component identities & 393\\
Random coupled systems, included in the preceding category & 16\\
Random finite-prefix versus resonance slope comparisons & 16\\
Geometric and cascade closed-form coefficient comparisons & 22\\
Finite prefixes of the accumulating support & 4\\
High-precision illustrations of the action-tail bound & 12\\
\bottomrule
\end{tabular}
\end{center}

The block identities cover Jordan sizes one through four, five
nonresonant/resonant gaps (including negative and small rational
gaps), and four polynomial degree ranges. Resonant constant terms
are checked explicitly. The scalar geometric example is compared
through action sixteen; the cascade is checked with six components.
The algebraic cancellation example and the Jordan/nonlinear example
are checked symbolically. Sixteen random coupled systems are computed
through action eight. A forcing perturbation at action five leaves
their earlier resonance certificates unchanged, as predicted.

The accumulating example is checked on generator prefixes ending at
\(n=3,6,12,24\), through action two, with symbolic
\(a,b,c,k\). These are finite prefixes of finite-generator models;
they are not a computational proof about the whole infinite monoid.
The supplied recurrence is safe only on an action set that is a
complete prefix for the finite input, or is divisor closed for the
requested targets. Missing an ancestor action invalidates an
arbitrary truncation. This precondition is documented in the code.

\subsection{A numerical illustration with an exact analytic bound}
For \eqref{eq:geometric-solution} with \(c=0\), take
\(r=0.05\), \(T=8\). Its coefficient norm is exactly
\[
 M_y=\sum_{n\ge1}(rT)^n=\frac{2}{3}.
\]
After retaining actions through \(N\), the exact remainder is
\(t^{N+1}/(1-t)\), where \(t=x\log x\). The theorem bounds its
absolute value by \((2/3)q_1(x)^{N+1}\).
A few recorded values are shown below; mantissas are rounded.

\begin{center}
\begin{tabular}{@{}rrrr@{}}
\toprule
\(x\) & \(N\) & Actual absolute remainder & Proved bound\\
\midrule
0.001 & 4 & \(1.56205\cdot10^{-11}\) & \(2.13333\cdot10^{-9}\)\\
0.005 & 4 & \(1.27111\cdot10^{-8}\) & \(6.66667\cdot10^{-6}\)\\
0.010 & 8 & \(8.90549\cdot10^{-13}\) & \(3.41333\cdot10^{-7}\)\\
0.030 & 12 & \(1.74820\cdot10^{-13}\) & \(8.70713\cdot10^{-4}\)\\
\bottomrule
\end{tabular}
\end{center}

The full JSON contains all twelve cases evaluated at 80-digit working
precision. These floating-point comparisons illustrate the proved
inequality; they are not interval-arithmetic certificates. The
conservative coefficient-norm estimate is not advertised as an
optimal pointwise bound.

\subsection{What is, and is not, formalized}
The inspected support module names the additive lemma
\path{Fabius.neumann_finite_decompositions} and the corresponding
support-preservation bridge \cite{repo-support}. These provide a
natural Lean starting point for Lemma~\ref{lem:divisor}. No new Lean
code is included in this package, and the general theorems have not
been checked by a proof assistant.

A feasible staged development would first formalize finite additive
divisors and divisor-closed recursion. The polynomial solver can then
be expressed using a nilpotent matrix, the polynomial derivative,
and a zero-constant integration operator. The degree budget and exact
slope theorem require only well-founded recursion, finite
coefficient convolutions, and degree inequalities. Algebraic
parameter dependence can be separated into a finite polynomial-map
construction. Finally, the norm estimates and analytic evaluation
can be layered on top of this algebraic core using a weighted
\(\ell^1\) space and the contraction mapping theorem.

The most important interface decision is to keep \emph{formal
support summability}, \emph{coefficient absolute convergence}, and
\emph{analytic evaluation} as different structures. The examples in
this article show that identifying any two without hypotheses would
lose precisely the distinction being studied.

\section{Further research questions and concrete next targets}
\label{sec:research}

The following questions are proposed extensions, not claims that each
is an established named open problem. Their value is that they have
specific hypotheses, potential obstructions, and test cases.

\subsection{Tail slope rather than maximal slope}
Theorem~\ref{thm:slope} computes the supremum over all actions. Determine
when the same value equals
\(\limsup_{\gamma\to\infty}\ell(P_\gamma)/\gamma\), with an
appropriate convention for finite support. The geometric example
propagates its resonant slope forever, whereas \(y=xL\) does not.
A useful first target is polynomial \(F\) on a finitely generated
monoid: characterize persistent maximal-slope propagation through
nonlinear dependency paths, including coefficient cancellations.
A graph-only criterion will generally need extra noncancellation
hypotheses.

\subsection{The full action--degree region}
The scalar \(\Theta\) forgets component information and most of the
support geometry. Determine the smallest convex region containing
all pairs \((\gamma,d)\) with nonzero coefficient, or a componentwise
version that records which nonlinear couplings are possible.
The degree-budget construction suggests a finite resonance-generated
upper region, but equality requires more than an upper-envelope
argument. A first task is to compare such regions in triangular
systems, where leading logarithmic coefficients can be computed
without solving all lower-degree coefficients.

\subsection{Several accumulation points and weighted divisibility}
Theorem~\ref{thm:sharp-family} works because all dangerous actions lie
below the first nontrivial quadratic interaction. For several
accumulation points, nonlinear sums may land near another resonance.
Find input-specific necessary and sufficient weighted coefficient
conditions that remain closed under convolution and the Euler
inverse. The first genuinely new test case should include two
accumulation sequences whose sums approach a third resonant action;
otherwise independent scalar weighting may conceal the main issue.

\subsection{Canonical summation when absolute realization fails}
For divergent normalized expansions, construct a canonical
renormalized solution class and determine its ambiguity. In the
linear repository example an infinite low-action block can be
renormalized; nonlinear feedback raises the question whether the
renormalization is compatible with products and composition.
A precise target is the family \eqref{eq:accum-ode} with
\(1<p\le2\): specify a subtraction prescription, prove that it
solves the original analytic equation, and identify the resulting
additional asymptotic scales rather than calling an arbitrary
regularization a Hahn sum.

\subsection{Higher logarithmic depth}
Replace logarithm-free coefficients by finite iterated-logarithmic
blocks and ask for an analogue of the exact resonance certificate.
The linear residue and logarithmic-depth infrastructure in ProveIt
is relevant, but nonlinear multiplication can propagate degrees
across more than one logarithmic level. The first target is a
triangular system with coefficients polynomial in \(\log\log(1/x)\)
and a positive outer-action gap. One must distinguish a new
logarithmic \emph{depth} from unbounded polynomial \emph{degree} at
an already present depth.

\subsection{Nonreal spectrum and oscillatory monomials}
The real-spectrum hypothesis lets actions and resonances live in
one ordered real monoid. With eigenvalues \(a+ib\), homogeneous
solutions introduce oscillatory factors \(e^{ib\Log x}\). Develop
an ordered support space that separates real growth from imaginary
frequency, and determine when finite divisors survive that
projection. A finite frequency lattice is a natural first case;
arbitrary accumulating frequencies may require a topology rather
than only a Hahn order. The analytic gap must then be formulated in
the complex spectral plane.

\subsection{Derivative-dependent nonlinearities and higher order}
Our source \(F(x,y)\) contains no derivatives of \(y\). For equations
involving \(Dy,D^2y,\ldots\), differentiation contributes factors
comparable to the action and may destroy the bounded nonlinear map
used in the contraction. Extend the finite-resonance maximum to
such differential polynomials, or construct a counterexample that
shows the correct invariant is different. The established
higher-order Dulac convergence theory provides a comparison point;
the non-locally-finite support and exact finite certificate remain
the specific targets here.

\subsection{Non-Archimedean actions}
In an ordered abelian group of higher rank, positive actions need
not have a real lower bound that controls the length of a dependency
path. The formal finite-antidiagonal lemma may survive while the
bound \(|k|\le\gamma/\delta\) has no Archimedean meaning.
Identify a positive grading into \(\RR\) that restores the proof,
and describe what happens when no such grading is available.
A lexicographic rank-two monoid should be the first test, with the
formal recursion and the analytic norm questions treated separately.

\subsection{Effective discovery of finite certificates}
Given a computable description of a support, when can one compute
\(\Div_\G(\lambda)\) and certify that it is complete? Theorem~\ref{thm:certificate}
is not a decision procedure for arbitrary real supports. For
restricted differential-algebraic inputs, compare the certificate
with the zero-testing framework of Chen, Fang, and van der Hoeven
\cite{chen-fang-vdh}, whose 2026 preprint concerns D-algebraic
transseries. Arbitrary coefficient families in the present article
need not be D-algebraic, so no automatic transfer of decidability is
implied. Exact finite-generator and algebraic-action presentations
are reasonable initial domains.

\subsection{Uniformity as resonances move}
Allow the eigenvalues and source parameters to vary. For fixed
matrix data, the slope sublevels are algebraic; as an eigenvalue
crosses a monoid action, the normalization changes and a logarithm
can appear. Construct a uniform local chart across such a crossing,
with explicit estimates that replace the blow-up of \(g^{-1}\).
An accumulation point of actions meeting a moving eigenvalue is a
more singular problem than one isolated crossing. The two should
not be treated by the same finite-dimensional normal form without
an additional argument.

\subsection{Certified infinite-block computation}
An action-tail estimate is useful computationally only after the
retained lower-action block is represented with a certified error.
Develop data structures for a finite divisor certificate together
with summable infinite blocks, and propagate explicit \(\ell^1\)
norm bounds through multiplication, integration, and nonlinear
fixed points. The low-action core \eqref{eq:low-core} is a concrete
benchmark. A successful implementation should report separately
its finite algebraic certificate, its block approximation error,
and its remaining action-tail bound.

\subsection{Formal verification of the finite-control theorem}
A focused Lean project could target the chain
\(\text{finite divisors}\to\text{normalized block recursion}
\to\text{exact slope}\to\text{algebraic strata}\)
before tackling analysis. The finite recursion must be stated in a
way that does not assume a locally finite action set. The
accumulating monoid with \(\Div_\G(2)=\{0,1,2\}\) is a useful
regression case. A later analytic layer should formalize the
restricted source space at resonances; omitting that degree margin
would make the claimed bounded right inverse false.

\section{Conclusion}
The nonlinear extension of finite logarithmic control is not a
finite global degree bound. It is the exact maximal degree-to-action
ratio, attained at a resonance and computed from finitely many
divisor coordinates. This invariant survives logarithm-free local
coordinate changes and yields algebraic loci of cancellations.
The underlying nonlinear solution may nevertheless contain
arbitrarily high logarithmic powers.

Convergence answers a different question. A uniform gap to the
spectrum is exactly what guarantees analytic realization for every
absolutely convergent source on a fixed real Hahn monoid. When the
gap fails, finite resonance data does not detect all analytic
obstructions. The Riccati family proves that this separation
persists under nonlinear feedback: the polynomial \(b+ca^2\)
controls logarithms, while \(\sum n|t_n|\) controls convergence.
These two certificates---one finite and algebraic, the other
potentially infinite and analytic---should remain distinct in both
symbolic implementations and future formalizations.

\appendix
\section{Hypothesis and dependency audit}
\label{app:audit}

The positive minimum \(\delta\) and real Archimedean action are used
to bound nonlinear arity, obtain a uniform valuation gain, and turn
a finite set of resonances into a linear degree budget. Well-ordering
is used both for recursion and for finite additive divisors. Neither
property may be replaced silently by local finiteness, and neither
is merely a convergence assumption.

The conditions \(f_{0,0}=f_{0,e_j}=0\) make the coefficient recurrence
strictly triangular in action. A zero-action linear term must first
be put in \(A\). A zero-action constant would leave the positive
ideal and is outside the theorem. The absence of logarithms in the
source coefficients is used in the exact slope proof: an explicit
\(L\)-factor at a nonresonant forcing action could create a new
maximal slope there, contradicting the stated formula.

Real spectrum is used to formulate resonant homogeneous modes and
small divisors in the same action order. Nilpotent Jordan parts are
retained throughout; diagonalizability is not assumed. The global
slope uses full vector degrees, so it is insensitive to a fixed
invertible choice of basis.

Formal existence and the finite certificate require no analytic
input condition and no gap. The analytic theorem requires the
absolute coefficient majorant \eqref{eq:input-norm}. The gap is
necessary for the \emph{universal} assertion only, and the sharp
family deliberately violates it while allowing many convergent
inputs. A convergent coefficient norm, not just well-ordered support,
is what permits all analytic sums and substitutions in the proof.

The theorem on coordinate invariance is restricted to invertible
constant linear parts and logarithm-free coefficients. It does not
claim invariance under singular shears, changes of the independent
variable, or transformations containing logarithms. The maximal
slope is not a tail limit, a pole order, or an exact radius of
analyticity.

Finally, the verification program checks finite exact identities and
illustrative numerical inequalities. It does not prove the infinite
support, Banach-space, or analytic-evaluation theorems. Those rely on
the conventional proofs in the article, and no proof-assistant
validation or independent peer review is claimed.

\Needspace{8\baselineskip}
\section{Repository provenance and inspection scope}
\label{app:provenance}

The observed repository snapshot used for this article is
\begin{center}\snap\end{center}
The connected GitHub tools were used to inspect the
transseries directory structure, the documentation index, the
Hahn--Fuchsian and residue-obstruction package descriptions, a
portion of the Hahn--Fuchsian source, and the complete support
module named in the bibliography. The large documentation index
and the Hahn--Fuchsian source exceeded the response display budget;
no complete line-by-line audit of those large files, the canonical
volume, or the whole repository is claimed. Newly delivered
archives at the repository's intake were not exhaustively audited.

The source-level support bridge was read at the pinned snapshot.
The mathematical theorems in the present article are proved
self-containedly rather than assumed from an unreviewed repository
report. External literature was consulted through primary author,
journal/preprint metadata, and paper text. The Dulac paper's PDF
text was accessible, but the remote screenshot endpoint failed;
no result here relies on an unseen figure or graphical measurement.
The PDF of this article was separately rendered and visually checked
in the local build workflow.

No repository files were modified. A suitable proposed placement is
\path{Analysis/Transseries/docs/series-and-transseries/Nonlinear_Hahn_Dulac/}.
This is a packaging suggestion, not an upload that has been performed.
The archive is standalone and does not load the repository's large
canonical volume or notation files.

\section{Reproduction instructions}
\label{app:reproduce}

The archive contains the editable \texttt{nonlinear\_hahn\_dulac.tex},
the compiled PDF, \texttt{build.sh}, the verification program and
recorded JSON, pinned verification requirements, and separate
proof/provenance/build notes. The bibliography is embedded in the
LaTeX source; no bibliography processor is needed.

To rebuild the article with a standard TeX Live installation, run
\begin{verbatim}
sh build.sh
\end{verbatim}
To repeat the finite verification, run
\begin{verbatim}
python -m pip install -r verification/requirements.txt
python verification/verify.py
\end{verbatim}
The verification program uses no network access after its dependencies
are installed. Random tests use a fixed seed. A successful run writes
\texttt{verification/results.json}. Exact algebraic assertions raise
an exception on failure; the program must be run without Python's
\texttt{-O} option, which would disable assertions.
% ed. (2026-09-29): output location changed on filing
\begin{ednote}
As filed, the program writes \path{build/results.json} (relative to the
package root) by default, so a rerun leaves the recorded
\path{verification/results.json} unchanged; writing the recorded file
requires \texttt{--overwrite-recorded}.
\end{ednote}

\clearpage
\begin{thebibliography}{9}
\raggedright

\bibitem{repo-support}
Vladimir Reshetnikov, \emph{ProveIt: well-based supports},
\path{Analysis/Transseries/Lean/Transseries/TransseriesWellBased.lean},
inspected at commit \snap, 29 September 2026.
\url{https://github.com/VladimirReshetnikov/ProveIt/blob/3d5973524506411392a911470b5ddc35521568ea/Analysis/Transseries/Lean/Transseries/TransseriesWellBased.lean}.

\bibitem{repo-residue}
\emph{Finite Residue Obstructions and Logarithmic-Depth Promotion in Hahn
Transseries}, research package in Vladimir Reshetnikov's ProveIt
repository, 29 September 2026; package README and editorial status notes.
\url{https://github.com/VladimirReshetnikov/ProveIt/tree/3d5973524506411392a911470b5ddc35521568ea/Analysis/Transseries/docs/series-and-transseries/Residue_Obstructions_Logarithmic_Depth_Promotion}.

\bibitem{repo-fuchsian}
\emph{Finite Resonance Certificates and Sharp Analytic Normalization for
Hahn--Fuchsian Systems}, research package in Vladimir Reshetnikov's
ProveIt repository, 29 September 2026; package README and inspected
source portion.
\url{https://github.com/VladimirReshetnikov/ProveIt/tree/3d5973524506411392a911470b5ddc35521568ea/Analysis/Transseries/docs/series-and-transseries/Hahn_Fuchsian_Resonance_Analytic_Normalization}.

\bibitem{vdh}
Joris van der Hoeven, \emph{Operators on generalized power series},
Illinois Journal of Mathematics \textbf{45}(4) (2001), 1161--1190.
\url{https://doi.org/10.1215/ijm/1258138061}.
Author's full-text version:
\url{https://www.texmacs.org/joris/noeth/noeth.html}.

\bibitem{gg-generalized}
Renat Gontsov and Irina Goryuchkina,
\emph{On the convergence of generalized power series satisfying an
algebraic ODE}, Asymptotic Analysis \textbf{93}(4) (2015), 311--325.
Preprint arXiv:1407.1330 (2014).
\url{https://arxiv.org/abs/1407.1330}.

\bibitem{gg-dulac}
Irina Goryuchkina and Renat Gontsov,
\emph{On the convergence of formal Dulac series satisfying an algebraic
ODE}, Sbornik: Mathematics \textbf{210}(9) (2019), 1207--1221.
Preprint arXiv:1712.08044 (2017), especially Lemmas 2--3.
\url{https://arxiv.org/abs/1712.08044}.

\bibitem{chen-fang-vdh}
Shaoshi Chen, Hanqian Fang, and Joris van der Hoeven,
\emph{A zero-test for D-algebraic transseries},
arXiv:2602.01188, submitted 1 February 2026.
\url{https://arxiv.org/abs/2602.01188}.

\end{thebibliography}
\end{document}
